\chapter{الآلة كاملة}
% full version: chapters/10_synthesis.tex

\pencilfig{10}{\pencilart{10}}{الآلة كاملة: ناقل البيانات، والمكابح فيه، وأربع قنوات تحكم إذاعية في الأعلى.}

كنا نأخذ نوعًا واحدًا من العصبونات في كل مرة ونسأل ماذا يضيف إلى الآلة. حان وقت جمع القطع معًا.

خرجت لنا آلة من ثلاث طبقات. في الأسفل شبكة هاتفية هائلة من عصبونات \nt{GLUT}: تجري فيها البيانات. وبجوارها العصبونات الكابحة \nt{GABA}: تتحكم في التدفق، وتقرر ما يمر، ومتى، وبأي علوّ. وفوق ذلك كله بضع محطات إذاعية، لكل منها مقبضها. \nt{DOPA} تقول أي تغييرات الوصلات تُستبقى. و\nt{SERO} تحفظ المستوى العام، ووفق الفرضية أفق الصبر. و\nt{NORA} تختار بين البحث والحسم. و\nt{ACET} بين الكتابة والقراءة. كل محطة هي في الحقيقة حزمة صغيرة من القنوات لا سلك واحد، لكن القنوات تبقى حفنة مقابل ستة وثمانين مليار عصبون.

\section*{قصة واحدة}

لنتتبع حدثًا واحدًا عبر الآلة كلها. يعرف الحيوان منذ زمن: اضغط الذراع أ تنل عصيرًا. وذات يوم يتغير العالم: أ لم يعد يعمل، والعصير يأتي به الذراع ب، الذي لم يكن الحيوان يستخدمه تقريبًا.

لم يصل العصير بعد أ، فيستجيب الدوبامين بفجوة، وتضعف الوصلات التي اختارت أ. تتكرر الفجوات، وتصير الأخطاء أكبر من المعتاد، فيخبر النورأدرينالين، بحسب الفرضية: العالم تغير. ينخفض مقبض التباين، ويلين الاختيار، ويدفع الضجيج أكثر إلى تجربة ب. وينقل الأسيتيل كولين الشبكة إلى وضع الكتابة. تجربة ب تجلب عصيرًا غير متوقع، فيطلق الدوبامين رشقة، وتقوى الوصلات التي اختارت ب. وبعد بضع محاولات يلحق التوقع بالعالم الجديد، وتنتهي المفاجآت، وتعود المقابض إلى مواضعها.

لاحظ: لا مقبض يعرف ما تفعله المقابض الأخرى. كل منها يستجيب لعلاماته، ويتكوّن ترتيب الأفعال وحده، كما في دائرة غير متزامنة يعمل فيها كل جزء حين تجهز مداخله. أن كل مقبض يتصرف هكذا بمفرده أمر مقيس في معظمه. أما أنها تعمل معًا بتناسق فلا يزال فرضية.

\section*{بمَ تختلف عن حاسوبك المحمول}

للحاسوب المحمول ساعة مشتركة، وللدماغ لا. عناصر الحاسوب موثوقة، والمشبك يضيّع أكثر من نصف الإشارات. ذاكرة الحاسوب منفصلة عن المعالج، وفي الدماغ هي في الوصلات نفسها. الحاسوب يتعلم (إن تعلم) بالانتشار العكسي للخطأ، والدماغ بقواعد محلية وبإشارة بث واحدة: ”أفضل مما كان متوقعًا“. والأهم: الحاسوب ينفق عشرات الواطات على معالج واحد، والدماغ عشرين واطًا على كل شيء.

\section*{ما لا نعرفه}

لا نعرف إلى أي حد تستخدم القشرة التوقيت الدقيق للنقرات. ولا نعرف كيف تضبط مليارات الوصلات في دوائر متعددة الطبقات، إذا كان معلمها واحدًا للجميع. ولا نفهم تمامًا ما الذي تنقله المحطات الإذاعية. ولا نعرف كيف تنسق أجزاء الدماغ المتباعدة عملها بلا ساعة مشتركة. ولا أحد حتى الآن يعرف كيف يبني آلة تفعل الشيء نفسه بعشرين واطًا.

بعد ذلك يخرج الكتاب عن حدود هذه النواة: إلى مداخل الآلة ومخارجها، وإلى تصحيح أعطالها، وإلى النوم، وإلى آلات من عصبونات حية على رقاقة.

\fullversion{10}
